% Generic section skeleton for statistical manuscripts.

\section{<SECTION TITLE>}\label{sec:<slug>}

% Section objective (1 sentence):
% Explain what this section accomplishes and why it is needed.

\subsection{<Subsection 1>}\label{sec:<slug>-part1}

% Paragraph 1: topic sentence + context.
% Paragraph 2: key content (method/data/result depending on section type).
% Paragraph 3: interpretation and transition.

\subsection{<Subsection 2>}\label{sec:<slug>-part2}

% Keep notation consistent and define symbols on first use.
% If details are missing, use placeholders such as \todo{add detail}.

% Cross-reference examples:
% See Section~\ref{sec:<related-section>}.
% As shown in Figure~\ref{fig:<figure-label>} and Table~\ref{tab:<table-label>}.
% From Equation~\eqref{eq:<equation-label>} we obtain ...

% Optional figure skeleton (prefer vector graphics .pdf/.eps for line plots):
% \begin{figure}[tbp]
%   \centering
%   % \includegraphics[width=0.85\textwidth]{figures/<file>.pdf}
%   \caption{<Self-contained caption>.}
%   \label{fig:<figure-label>}
% \end{figure}

% Optional table skeleton:
% \begin{table}[tbp]
%   \centering
%   \caption{<Self-contained caption>.}
%   \label{tab:<table-label>}
%   \begin{tabular}{...}
%   \end{tabular}
% \end{table}
